\documentclass[11pt,professionalfont,aspectratio=43]{beamer}

% Assets
\usepackage[czech]{babel}				% Jazyk
\usepackage[a-2u]{pdfx}				    % Kopírování z pdfka
\usepackage{tikz}						% Schémata automatů
\usetikzlibrary{shadows, shadows.blur, shapes.geometric, arrows, positioning, overlay-beamer-styles, patterns, shadings, shapes.callouts}
% \usetikzlibrary{arrows}
\usepackage{pgfplots}
% \pgfplotsset{compat=1.15}
% \usepackage{mathrsfs}
\usepackage[export]{adjustbox}

\usepackage[T1]{fontenc}
\usefonttheme[onlymath]{serif}
% \usepackage{newtxmath}
\usepackage{pxfonts}
\usepackage{slantsc}
% \usepackage{varwidth}

\usepackage[linesnumbered,ruled,vlined]{algorithm2e}                % Pseudokód algoritmů
\usepackage[utf8]{inputenc}
% \usepackage{textcomp}
\usepackage{hyperref}
\usepackage{fancybox}

%%% FONTS
% \usefonttheme{serif}
% \usepackage{lmodern}
\usepackage{helvet}

\hypersetup{%
    pdfencoding=auto,
    pdfauthor={\insertauthor},
    pdftitle={\insertsubtitle}
}
\usepackage{float}
\usepackage{csquotes}					% české uvozovky
\usepackage{enumerate}					% enumerate environment
\usepackage{indentfirst}
\usepackage{tcolorbox}                  % Fancy boxíky
\usepackage{mathtools}
\usepackage{pifont}
\usepackage{cancel}
\usepackage{subcaption}
\usepackage{soul}
\usepackage{makecell}                   % Zalomení v buňce tabulky
\usepackage{xcolor}
\usepackage{graphicx}
\usepackage{tabularx}
\usepackage{amsmath}
\usepackage{amsthm}
\usepackage{colortbl}                   % Barvení buněk v tabulce
\usepackage{multicol}                   % Rozdělení obsahu do sloucpů

\usepackage{listings}                   % Úryvky z kódu
\definecolor{maroon}{rgb}{0.5,0,0}
\definecolor{forestgreen}{rgb}{0.13, 0.55, 0.13}

% \lstdefinelanguage{XML}
% {
%   basicstyle=\ttfamily,
%   morestring=[s]{"}{"},
%   morecomment=[s]{?}{?},
%   morecomment=[s]{!--}{--},
%   commentstyle=\color{forestgreen},
%   moredelim=[s][\color{black}]{>}{<},
%   moredelim=[s][\color{red}]{\ }{=},
%   stringstyle=\color{blue},
%   identifierstyle=\color{maroon}
% }
% Colors

\def\maincolR{0} % červená složka (0–255)
\def\maincolG{83} % zelená složka
\def\maincolB{120} % modrá složka
\definecolor{maincolor}{HTML}{8B00B5}

% \pgfmathtruncatemacro{\authorfooterbgR}{1*\maincolR}
% \pgfmathtruncatemacro{\authorfooterbgG}{1*\maincolG}
% \pgfmathtruncatemacro{\authorfooterbgB}{1*\maincolB}
% \definecolor{authorfooterbg}{RGB}{\authorfooterbgR,\authorfooterbgG,\authorfooterbgB}

% \pgfmathtruncatemacro{\titlefooterbgR}{0.722*\maincolR}   
% \pgfmathtruncatemacro{\titlefooterbgG}{0.974*\maincolG}   
% \pgfmathtruncatemacro{\titlefooterbgB}{0.739*\maincolB}
% \definecolor{titlefooterbg}{RGB}{\titlefooterbgR,\titlefooterbgG,\titlefooterbgB}

% Beamer theme
\colorlet{themecolor}{maincolor}
\colorlet{authorfooterbg}{maincolor}
\definecolor{titlefooterbg}{HTML}{9471B0}
\definecolor{titlefooterfg}{HTML}{000000}
\colorlet{titlecolorfg}{maincolor}
\definecolor{titleboxbg}{HTML}{E7E7E7}
\definecolor{datefooterbg}{HTML}{E3E3E3}
\colorlet{titleframeboxcolor}{maincolor}

% Background gradient
\definecolor{gradMiddle}{RGB}{212,248,255}
\definecolor{gradEnd}{RGB}{114,231,255}

% Block theme
\definecolor{blocktitlebg}{HTML}{F0DFA9}
\definecolor{blocktitlefg}{HTML}{957300}
\definecolor{blockbodybg}{HTML}{F5F0DF}

\definecolor{alertblockbodybg}{HTML}{FFE4E4}
\definecolor{alertblocktitlebg}{HTML}{ECA4A4}
\definecolor{alertblocktitlefg}{HTML}{B11313}

\definecolor{exampleblockbodybg}{HTML}{DAF0DA}
\definecolor{exampleblocktitlebg}{HTML}{BBE0BB}
\definecolor{exampleblocktitlefg}{HTML}{257A25}

\definecolor{mathfg}{HTML}{00185E}

\colorlet{bulletbg}{maincolor}
% \definecolor{bulletbg}{HTML}{}

% Beamer theme
\usetheme{Madrid}
\usecolortheme{seahorse}
% \useoutertheme[subsection=false]{miniframes}

\usecolortheme[named=themecolor]{structure}
\setbeamertemplate{frame numbering}[fraction]
\setbeamertemplate{navigation symbols}{}
\setbeamercovered{invisible}

% Header/footer
\setbeamercolor{author in head/foot}{bg=authorfooterbg,fg=white}
\setbeamercolor{title in head/foot}{bg=titlefooterbg,fg=titlefooterfg}
\setbeamercolor{date in head/foot}{bg=datefooterbg,fg=black}
% \setbeamercolor{title}{bg=datefooterbg}
% \setbeamercolor{title}{fg=titlecolorfg}
\setbeamerfont{title}{series=\bfseries}
\setbeamercolor{frametitle}{bg=datefooterbg}
% \setbeamercolor{section in head/foot}{bg=red,fg=cyan}
% \setbeamercolor{subsection in head/foot}{bg=blue,fg=yellow}

% Table of contents
\setbeamercolor{section in toc}{fg=black}
\setbeamercolor{subsection in toc}{fg=red}
\setbeamercolor{section number projected}{fg=black}

% Blocks
\setbeamertemplate{blocks}[rounded][shadow=true]

\setbeamercolor{block title}{bg=blocktitlebg, fg=blocktitlefg}
\setbeamercolor{block body}{parent=normal text,use=block title,bg=blockbodybg}

\setbeamercolor{block body alerted}{bg=alertblockbodybg}
\setbeamercolor{block title alerted}{bg=alertblocktitlebg,fg=alertblocktitlefg}

\setbeamercolor{block body example}{bg=exampleblockbodybg}
\setbeamercolor{block title example}{bg=exampleblocktitlebg,fg=exampleblocktitlefg}

% Math
\setbeamercolor{math text}{fg=mathfg}

% Itemize a enumerate
\setbeamertemplate{itemize items}[circle]
\setbeamertemplate{enumerate items}[circle]
\setbeamercolor{itemize item}{fg=bulletbg,bg=white}
\setbeamercolor{itemize subitem}{fg=bulletbg,bg=white}
\setbeamercolor{itemize subsubitem}{fg=bulletbg,bg=white}

\setbeamertemplate{enumerate items}[circle]
\setbeamercolor{item projected}{bg=bulletbg}

% Titlepage
\defbeamertemplate*{title page}{customized}[1][]
{
    \vbox{}%
    \vfill%
    \begin{centering}
        \begin{tikzpicture}
            \node[thick, rectangle, rounded corners=5mm, drop shadow={shadow xshift=0mm, shadow yshift=2mm, opacity=0.2, blur shadow, shadow blur steps=10}, inner sep=5mm, fill=titleboxbg, fill opacity=1] {%
                {\usebeamerfont{title}\textcolor{titlecolorfg}{\inserttitle}\par}%
            };
        \end{tikzpicture}
        \vskip1em\par%
        \ifx\insertsubtitle\@empty%
        \else%
            {\usebeamerfont{subtitle}\insertsubtitle\par}%
        \fi%
        \vskip1em\par%
        \ifx\insertauthor\@empty%
        \else%
            {\usebeamerfont{author}\insertauthor}\\%
            \texttt{\email}
        \fi
        \vskip0.3em\par%
        % \ifx\insertinstitute\@empty%
        % \else%
        %     {\usebeamerfont{institute}\insertinstitute\par}%
        % \fi%
        \vskip1em\par%
        \ifx\inserttitlegraphic\@empty%
        \else%
            \vskip1em\par%
            \inserttitlegraphic\par%
        \fi%
        \vskip1em\par%
        \ifx\insertdate\@empty%
        \else%
            {\usebeamerfont{date}\insertdate\par}%
        \fi%
    \end{centering}%
    \vfill%
}
% Enumerate
%\setlist[enumerate]{topsep=0pt,itemsep=-1ex,partopsep=1ex,parsep=1ex,label=(\arabic*)}

\MakeOuterQuote{"}

%%% Makra pro definice, věty, tvrzení, příklady, ... (vyžaduje baliček amsthm)

\theoremstyle{plain}
\newtheorem{thm}{Věta}[section]
\newtheorem{lem}[theorem]{Lemma}
\newtheorem{prop}[theorem]{Tvrzení}
\newtheorem{cor}[theorem]{Důsledek}
\newtheorem*{prop*}{Tvrzení}

\theoremstyle{definition}
\newtheorem{define}[theorem]{Definice}
% \newtheorem{example}[theorem]{Příklad}
% \newtheorem{remark}[theorem]{Poznámka}
% \newtheorem{convention}[theorem]{Úmluva}
% \newtheorem{denoting}[theorem]{Značení}

\newenvironment{defblock}[1][]{
  \begin{block}{\ifthenelse{\equal{#1}{}}
  {Definice}
  {Definice (#1)}}
}{\end{block}}
\newenvironment{thmblock}[1][]{
  \begin{block}{\ifthenelse{\equal{#1}{}}
  {Věta}
  {Věta (#1)}}
}{\end{block}}
\newenvironment{corblock}[1][]{
  \begin{block}{\ifthenelse{\equal{#1}{}}
  {Důsledek}
  {Důsledek (#1)}}
}{\end{block}}
\newenvironment{propblock}[1][]{
  \begin{block}{\ifthenelse{\equal{#1}{}}
  {Tvrzení}
  {Tvrzení (#1)}}
}{\end{block}}
\newenvironment{exblock}[1][]{
  \begin{exampleblock}{\ifthenelse{\equal{#1}{}}
  {Příklad}
  {Příklad (#1)}}
}{\end{exampleblock}}


% Colors
\definecolor{lightblue}{HTML}{009AD4}
\definecolor{lightgreen}{HTML}{68FF00}
\definecolor{darkgreen}{HTML}{00B600}
\definecolor{darkblue}{HTML}{0265b3}
\definecolor{darkred}{HTML}{DA0707}
\definecolor{lightred}{HTML}{FF5100}
\definecolor{orange}{HTML}{FFE000}
\definecolor{codeblue}{HTML}{007eed}
\definecolor{codegreen}{rgb}{0,0.6,0}
\definecolor{codegray}{rgb}{0.5,0.5,0.5}
\definecolor{codebeige}{HTML}{D4A000}
\definecolor{codepink}{HTML}{FF0055}
\definecolor{backcolour}{rgb}{0.95,0.95,0.92}

\definecolor{decproblemtitlefg}{HTML}{957300}
\definecolor{decproblembg}{HTML}{F5F0DF}

\definecolor{timportantboxcolor}{HTML}{A97100}
\definecolor{talertboxcolor}{HTML}{B70000}
\definecolor{tnoticeboxcolor}{HTML}{159C00}
\definecolor{tinfoboxcolor}{HTML}{004DFF}

\newcommand{\markred}[1]{\textcolor{lightred}{#1}}
\newcommand{\markdarkred}[1]{\textcolor{darkred}{#1}}
\newcommand{\markgreen}[1]{\textcolor{lightgreen}{#1}}
\newcommand{\markdarkgreen}[1]{\textcolor{darkgreen}{#1}}
\newcommand{\markorange}[1]{\textcolor{orange}{#1}}
\newcommand{\markblue}[1]{\textcolor{lightblue}{#1}}
\newcommand{\markdarkblue}[1]{\textcolor{darkblue}{#1}}

% Math mode settings
% \AtBeginDocument{
%     \setlength{\abovedisplayskip}{3pt}
%     \setlength{\belowdisplayskip}{-4pt}
% }

% Inline images
\newcommand{\inlineimgscale}{1.1}

% X and check mark
\newcommand{\cmark}{\markgreen{\ding{51}}}
\newcommand{\xmark}{\markred{\ding{55}}}

% Math
\newcommand{\R}{\mathbb{R}}
\newcommand{\C}{\mathbb{C}}
\newcommand{\N}{\mathbb{N}}
\newcommand{\Z}{\mathbb{Z}}
\newcommand{\Q}{\mathbb{Q}}

\newcommand{\T}[1]{#1^\top}                     % Tranpozice matice
\newcommand{\compl}[1]{#1^\complement}
\newcommand{\mapping}[3]{#1:#2 \rightarrow #3}
\newcommand{\mapto}[3]{#1:#2 \mapsto #3}
\newcommand{\set}[1]{\left\{#1\right\}}
\newcommand{\powset}{\mathcal{P}}
\newcommand{\code}[1]{\langle#1\rangle}
\DeclareMathOperator*{\argmin}{arg\,min}
\DeclareMathOperator*{\argmax}{arg\,max}

% Statistika a AI
\newcommand{\average}[1]{\overline{#1}}
\DeclareMathOperator{\median}{med}
\DeclareMathOperator{\Var}{var}
\DeclareMathOperator{\Cov}{cov}
\DeclareMathOperator{\MSE}{MSE}
\DeclareMathOperator{\SSE}{SSE}

\DeclareMathOperator{\TP}{TP}
\DeclareMathOperator{\TN}{TN}
\DeclareMathOperator{\FP}{FP}
\DeclareMathOperator{\FN}{FN}
\DeclareMathOperator{\Accuracy}{Acc}
\DeclareMathOperator{\Precision}{Prec}
\DeclareMathOperator{\Recall}{Rec}
\DeclareMathOperator{\Fscore}{F}
\DeclareMathOperator{\Gini}{Gini}

\DeclareMathOperator{\mspan}{span}
\DeclareMathOperator{\mrank}{rank}
\DeclareMathOperator{\tg}{tg}
\DeclareMathOperator{\id}{id}
\DeclareMathOperator{\dom}{dom}
\DeclareMathOperator{\rng}{rng}
\renewcommand{\emptyset}{\varnothing}
\newcommand{\binary}[1]{\ensuremath{\left\langle #1\right\rangle}_B}

% DFS in and out lists
\DeclareMathOperator{\inlist}{in}
\DeclareMathOperator{\outlist}{out}
% \DeclareMathOperator{\lowlist}{low}
\DeclareMathOperator{\key}{key}

% A* macro
\newcommand{\Astar}{$\textcolor{black}{\textup{A}^{*}}$}

% Complexity classes
\DeclareMathOperator{\classP}{P}
\DeclareMathOperator{\classNP}{NP}
\DeclareMathOperator{\classTime}{TIME}
\DeclareMathOperator{\classNTime}{NTIME}
\DeclareMathOperator{\classSpace}{SPACE}
\DeclareMathOperator{\classNSpace}{NSPACE}

\DeclareMathOperator{\SAT}{SAT}
\DeclareMathOperator{\THREESAT}{3-SAT}
\DeclareMathOperator{\UNSAT}{UNSAT}
\DeclareMathOperator{\NZMNA}{NzMna}
\DeclareMathOperator{\KLIKA}{Klika}
\DeclareMathOperator{\ROZVRH}{Rozvrh}
\DeclareMathOperator{\SUDOKU}{Sudoku}
\newcommand{\HALT}{\ensuremath{\textsc{Halt}}}
\newcommand{\FACTOR}{\ensuremath{\textsc{Factor}}}

\DeclareMathOperator{\regex}{RegE}

%%% Barevné boxíky (notice / info / important / alert)
%
% Šířka boxu se nezadává ručně -- box se přizpůsobí šířce svého obsahu.
% Obsah se vysází do prostředí varwidth, takže krátký text zůstane na
% jednom řádku; delší text se zalomí, nejvýše však na šířku \linewidth
% (tj. \textwidth mimo sloupce a minipage). Pokud je titulek širší než
% obsah, řídí šířku boxu titulek.

\newsavebox{\tboxcontentbox}
\newsavebox{\tboxtitlebox}
\newlength{\tboxwidth}
\newlength{\tboxpadding}

% Geometrie boxu (musí odpovídat stylu tautobox níže):
% 2 * (boxrule + boxsep + left)
\setlength{\tboxpadding}{\dimexpr 0.5mm+1mm+4mm\relax}
\setlength{\tboxpadding}{2\tboxpadding}

\tcbset{tautobox/.style={boxrule=0.5mm, boxsep=1mm, left=4mm, right=4mm,
                         fonttitle=\bfseries}}

% Začátek sběru obsahu do boxu
\newcommand{\tboxcollect}{%
  \begin{lrbox}{\tboxcontentbox}%
    \begin{varwidth}{\dimexpr\linewidth-\tboxpadding\relax}%
}

% Ukončení sběru a výpočet výsledné šířky
% #1 = titulek (pro měření; prázdný u variant bez titulku)
\newcommand{\tboxmeasure}[1]{%
    \end{varwidth}%
  \end{lrbox}%
  \sbox{\tboxtitlebox}{\bfseries #1}%
  \setlength{\tboxwidth}{\wd\tboxcontentbox}%
  \ifdim\wd\tboxtitlebox>\tboxwidth
    \setlength{\tboxwidth}{\wd\tboxtitlebox}%
  \fi
  \addtolength{\tboxwidth}{\tboxpadding}%
  \ifdim\tboxwidth>\linewidth
    \setlength{\tboxwidth}{\linewidth}%
  \fi
}

% Vysázení boxu s titulkem: #1 = barva, #2 = titulek
\newcommand{\tboxrenderbox}[2]{%
  \tboxmeasure{#2}%
  \begin{center}%
    \begin{tcolorbox}[tautobox, colback=#1!15, colframe=#1,
                      title={#2}, width=\tboxwidth]
      \usebox{\tboxcontentbox}%
    \end{tcolorbox}%
  \end{center}%
}

% Vysázení boxu bez titulku: #1 = barva
\newcommand{\tboxrenderframe}[1]{%
  \tboxmeasure{}%
  \begin{center}%
    \begin{tcolorbox}[tautobox, colback=#1!15, colframe=#1,
                      width=\tboxwidth]
      \usebox{\tboxcontentbox}%
    \end{tcolorbox}%
  \end{center}%
}

% Poznámky (notice)
\newenvironment{tnoticebox}[1]{%
  \def\tboxtitletext{#1}\tboxcollect
}{%
  \tboxrenderbox{tnoticeboxcolor}{\tboxtitletext}%
}

\newenvironment{tnoticeframe}{%
  \tboxcollect
}{%
  \tboxrenderframe{tnoticeboxcolor}%
}

% Informační bloky (info)
\newenvironment{tinfobox}[1]{%
  \def\tboxtitletext{#1}\tboxcollect
}{%
  \tboxrenderbox{tinfoboxcolor}{\tboxtitletext}%
}

\newenvironment{tinfoframe}{%
  \tboxcollect
}{%
  \tboxrenderframe{tinfoboxcolor}%
}

% Důležité bloky (important)
\newenvironment{timportantbox}[1]{%
  \def\tboxtitletext{#1}\tboxcollect
}{%
  \tboxrenderbox{timportantboxcolor}{\tboxtitletext}%
}

\newenvironment{timportantframe}{%
  \tboxcollect
}{%
  \tboxrenderframe{timportantboxcolor}%
}

% Varovné bloky (alert)
\newenvironment{talertbox}[1]{%
  \def\tboxtitletext{#1}\tboxcollect
}{%
  \tboxrenderbox{talertboxcolor}{\tboxtitletext}%
}

\newenvironment{talertframe}{%
  \tboxcollect
}{%
  \tboxrenderframe{talertboxcolor}%
}

% Decision problem
\newcommand{\decproblem}[3]{
    \begin{center}
        \begin{minipage}{.9\textwidth}
            \begin{table}[!htbp]
                \centering
                \bgroup
                \def\arraystretch{1.8}
                \begin{tabularx}{\textwidth}{|rX|}
                \hline
                \rowcolor{decproblembg} \multicolumn{2}{|c|}{\textcolor{decproblemtitlefg}{\textsc{#1}}} \\ \hline
                \markdarkred{\textbf{Vstup:}} & #2                \\ \hline
                \rowcolor{decproblembg} \markdarkred{\textbf{Otázka:}} & #3               \\ \hline
                \end{tabularx}
                \egroup
            \end{table}
        \end{minipage}
    \end{center}
}
\newcommand{\bigO}{\mathcal{O}}

% TODO
\newcommand{\todo}[1]{\textcolor{red}{(\noindent TODO: #1.)}}


\newcommand{\hlcode}[1]{\colorbox{red}{#1}}

% Algorithms
\numberwithin{algocf}{section}
\renewcommand{\algorithmcfname}{Algoritmus}
\SetKwInput{KwIn}{Vstup}
\SetKwInput{KwOut}{Výstup}
\SetKwProg{Function}{Funkce}{}{end}
\setlength{\algomargin}{25pt}
\renewcommand{\thealgocf}{}     % Vypnutí číslování algoritmů

% Definice stylů pro uzly
\tikzstyle{startstop} = [rectangle, rounded corners, minimum width=1.5cm, minimum height=0.5cm,text centered, draw=black, fill=red!30]
\tikzstyle{process}   = [rectangle, minimum width=1.5cm, minimum height=0.5cm, text centered, draw=black, fill=orange!30]
\tikzstyle{decision}  = [diamond, aspect=2, text centered, draw=black, fill=green!30]
\tikzstyle{arrow}     = [thick,->,>=stealth]

\tikzstyle{vertex}    = [thick, draw, circle, minimum size=4mm, inner sep=1pt, outer sep=1pt, align=center]
\tikzstyle{vertexplain}    = [minimum size=4mm, inner sep=1pt, outer sep=1pt, align=center]
\tikzstyle{verteximg}    = [draw, circle, minimum size=4mm, inner sep=0pt, outer sep=1pt, align=center, clip]
\tikzstyle{verteximgplain}    = [minimum size=4mm, inner sep=0pt, outer sep=1pt, align=center]
\tikzstyle{task}    = [draw, rounded corners=2pt, minimum width=14mm, minimum height=8mm, inner sep=2pt, align=center]

\tikzstyle{edge}    = [thick,>=stealth]
\tikzstyle{diredge}    = [thick, ->, >=stealth]

%% Automaty
\tikzstyle{state}=[%
    draw,
    circle,
    thick,
    minimum size=10mm,
    inner sep=0pt,
    outer sep=0pt
]

\tikzstyle{acceptstate}=[%
    state,
    double,
    double distance=1pt
]

\tikzstyle{transition}=[%
    ->,
    >=stealth,
    thick
]

\newcommand<>{\state}[4][]{%
    \node#5[state,#1] (#2) at #3 {#4};
}

\newcommand<>{\acceptingstate}[4][]{%
    \node#5[acceptstate,#1] (#2) at #3 {#4};
}

\newcommand<>{\transitionarrow}[5][]{%
    \path#6 (#2) edge[transition,#1] node[#3] {#4} (#5);
}

\newcommand<>{\initialstate}[6][]{%
    \node#7[state,#1] (#2) at #3 {#4};
    \draw#7[transition] #5 -- (#2.#6);
}

%%%%%

% Frames
\newcommand{\tableofcontentsframe}{
    \begin{frame}{Obsah}
        \tableofcontents
    \end{frame}
}

% \definecolor{boxcolor}{HTML}{00820a}
\newcommand{\titleframe}[1]{%
  \begin{frame}[plain]
    \begin{tikzpicture}[remember picture, overlay]
      \node at (current page.center)
        [draw=titleframeboxcolor, line width=1pt,
         inner xsep=10mm, inner ysep=4mm, rounded corners=2mm,
         text width=0.8\textwidth, align=center,
         execute at begin node=\setlength{\baselineskip}{2.2em}]
        {\Huge\markred{#1}\normalfont};
    \end{tikzpicture}
  \end{frame}
  \addtocounter{framenumber}{-1}%
}

\newenvironment{titleframeenv}[1]{
    \begin{frame}[plain]
        \begin{tcolorbox}[colback=white,boxrule=1pt]
            \begin{center}
                \Huge\markred{#1}\normalfont
            \end{center}
        \end{tcolorbox}
}{\addtocounter{framenumber}{-1}\end{frame}}

\newcommand{\icon}[1]{
  \includegraphics[height=\fontcharht\font`\B]{#1}
}

% Obrázek, který se vždy vejde na snímek. Adjustbox jej v případě potřeby
% zmenší tak, aby nepřesáhl šířku textu ani zadanou část výšky snímku; díky
% tomu sedí sazba ve všech používaných poměrech stran (4:3, 16:9 i 16:10).
% Volitelným argumentem je maximální výška (výchozí 0.68\textheight), povinným
% cesta k souboru s obrázkem (TikZ zdroj vkládaný přes \input).
\newcommand{\obrazek}[2][0.68\textheight]{%
    \begin{figure}
        \centering
        \begin{adjustbox}{max width=\linewidth, max totalheight=#1, center}
            \input{#2}
        \end{adjustbox}
    \end{figure}%
}

% Totéž pro rastrové či vektorové obrázky vkládané přes \includegraphics.
\newcommand{\obrazekgrf}[2][0.68\textheight]{%
    \begin{figure}
        \centering
        \includegraphics[width=\linewidth, max totalheight=#1, center]{#2}
    \end{figure}%
}

% Závěrečný snímek se seznamem zdrojů. Prvním (volitelným) argumentem je název
% sekce i snímku, druhým čárkou oddělený seznam klíčů z ../assets/literatura.bib.
% Citace se sázejí podle ČSN ISO 690, v pořadí, v jakém jsou klíče uvedeny.
\newcommand{\zdrojeframe}[2][Zdroje]{%
    \section{#1}
    \nocite{#2}%
    \begin{frame}[allowframebreaks]{#1}
        \printbibliography[heading=none]
    \end{frame}
}
\input{../assets/listing.tex}

\def\presentationtitle{Dynamické programování}

% Obrázek se vždy vejde na snímek: adjustbox jej v případě potřeby
% zmenší tak, aby nepřesáhl šířku textu ani zadanou část výšky snímku.
% Díky tomu sedí sazba ve všech používaných poměrech stran.
\newcommand{\obrazek}[2][0.68\textheight]{%
    \begin{figure}
        \centering
        \begin{adjustbox}{max width=\textwidth, max totalheight=#1, center}
            \input{#2}
        \end{adjustbox}
    \end{figure}%
}

% Ceník tyčí, se kterým se pracuje v celé kapitole o dělení tyče.
\newcommand{\tyccenik}{%
    \begin{center}
        \renewcommand{\arraystretch}{1.25}
        \begin{tabular}{|r|c|c|c|c|c|c|c|c|c|c|}
            \hline
            \rowcolor{decproblembg}
            délka $i$   & $1$ & $2$ & $3$ & $4$ & $5$ & $6$ & $7$ & $8$ & $9$ & $10$ \\ \hline
            cena $p_i$  & $1$ & $5$ & $8$ & $9$ & $10$ & $17$ & $17$ & $20$ & $24$ & $30$ \\ \hline
        \end{tabular}
    \end{center}
}

\begin{document}

    \title[\presentationtitle]{\textsc{Teoretická informatika}}
\subtitle{\presentationtitle}
\author{David Weber}
\def\email{weber3@spsejecna.cz}
\institute{SPŠE Ječná}
\date{%
    Rok 
    \pgfmathsetmacro{\firstyear}{ifthenelse(\month<9, \year-1, \year)}%
    \pgfmathsetmacro{\secondyear}{ifthenelse(\month<9, \year, \year+1)}%
    \pgfmathtruncatemacro{\secondyearlasttwo}{mod(\secondyear,100)}%
    \pgfmathprintnumber[fixed,precision=0,zerofill,1000 sep={}]{\firstyear}/%
    \pgfmathprintnumber[fixed,precision=0,zerofill,1000 sep={}]{\secondyearlasttwo}%
}
\titlegraphic{\includegraphics[width=0.25\textwidth]{../images/SPSE-Jecna_Logo.pdf}}

\begin{frame}[plain]
    \titlepage
\end{frame}

    % =================================================================
    \section{Opakování: rekurze}
    % =================================================================

    \begin{frame}{Co je rekurze?}
        \begin{itemize}
            \visible<2->{\item \markdarkred{Rekurzivní} je taková funkce, která ve svém těle \textbf{volá sama sebe}.}
            \visible<3->{\item Aby výpočet skončil, musí každá rekurze obsahovat:}
            \begin{itemize}
                \visible<4->{\item \textbf{triviální případ}, který umíme vyřešit přímo,}
                \visible<5->{\item \textbf{rekurzivní krok}, který úlohu převede na \textbf{menší} úlohu téhož druhu.}
            \end{itemize}
        \end{itemize}
        \visible<6->{\begin{tinfoframe}
            \begin{center}
                Rekurze je vlastně recept: \uv{umím-li vyřešit menší úlohu, umím vyřešit i tu svou}.
            \end{center}
        \end{tinfoframe}}
    \end{frame}

    \begin{frame}{Rekurze}{Příklad: faktoriál}
        \visible<2->{\[
            n!=
            \begin{cases}
                1              & \text{pro } n\leqslant 1,\\
                n\cdot(n-1)!   & \text{jinak.}
            \end{cases}
        \]}
        \visible<3->{\obrazek[0.55\textheight]{images/rekurze-faktorial.tex}}
    \end{frame}

    \begin{frame}{Rekurze počítá \uv{shora dolů}}
        \begin{itemize}
            \visible<2->{\item Rekurze začíná u \textbf{celého problému} a postupně jej rozkládá na menší a menší podproblémy.}
            \visible<3->{\item Teprve na dně, u triviálních případů, se začnou skládat výsledky zpět nahoru.}
            \visible<4->{\item Tomuto směru výpočtu říkáme \markdarkred{shora dolů} (angl. \emph{top-down}).}
        \end{itemize}
        \visible<5->{\begin{tnoticeframe}
            \begin{center}
                Rekurze \textbf{nijak neřeší}, zda tentýž podproblém neřešíme opakovaně.
            \end{center}
        \end{tnoticeframe}}
    \end{frame}

    % =================================================================
    \section{Dynamické programování}
    % =================================================================

    \titleframe{Dynamické programování}

    \subsection{Co to je?}

    \begin{frame}{Co je dynamické programování?}
        \begin{itemize}
            \visible<2->{\item Technika návrhu algoritmů, která stejně jako \textbf{Rozděl a panuj} staví na \markdarkred{rekurzivním rozkladu} problému na podproblémy.}
            \visible<3->{\item Navíc ale využívá toho, že se \textbf{tytéž podproblémy během rekurze opakují}.}
            \visible<4->{\item Místo opakovaného počítání si už jednou spočítané výsledky \markdarkred{pamatuje v tabulce}.}
        \end{itemize}
        \visible<5->{\begin{tinfobox}{Odkud ten název?}
            Termín zavedl Richard Bellman v~50. letech. \uv{Programováním} se tehdy mínilo \textbf{plánování} a \uv{dynamické} odkazuje na to, že optimální volba v~každém kroku závisí na krocích předchozích.
        \end{tinfobox}}
    \end{frame}

    \begin{frame}{Kdy dynamické programování funguje?}
        \visible<2->{\begin{defblock}[Optimální podstruktura]
            Úloha má \markdarkred{optimální podstrukturu}, pokud je optimální řešení celé úlohy složené z~optimálních řešení jejích podúloh.
        \end{defblock}}
        \visible<3->{\obrazek[0.36\textheight]{images/optimalni-podstruktura.tex}}
        \visible<4->{\begin{timportantframe}
            \begin{center}
                Bez této vlastnosti bychom z~vyřešených podúloh nedokázali poskládat řešení té celé.
            \end{center}
        \end{timportantframe}}
    \end{frame}

    \begin{frame}{Proč se to hodí?}
        \begin{itemize}
            \visible<2->{\item Každý podproblém vyřešíme \textbf{právě jednou} a výsledek si zapamatujeme.}
            \visible<3->{\item Tím se z~exponenciálně pomalého algoritmu často stane \markdarkred{polynomiální}.}
            \visible<4->{\item Platíme za to pamětí na tabulku $\implies$ jde o~výměnu \textbf{paměti za čas}.}
            \visible<5->{\item Pokud je počet různých podproblémů polynomiální a každý z~nich vyřešíme v~polynomiálním čase, běží celý algoritmus v~polynomiálním čase.}
        \end{itemize}
    \end{frame}

    \subsection{Rozdíl oproti rekurzi}

    \begin{frame}{Shora dolů vs. zdola nahoru}
        \begin{itemize}
            \visible<2->{\item \markdarkred{Rekurze} jde \textbf{shora dolů}: začíná u celého problému a ptá se na menší podproblémy.}
            \visible<3->{\item \markdarkred{Dynamické programování} jde \textbf{zdola nahoru}: začne u nejmenších podproblémů a postupuje k~tomu zadanému.}
        \end{itemize}
        \visible<4->{\obrazek[0.62\textheight]{images/shora-dolu-zdola-nahoru.tex}}
    \end{frame}

    \begin{frame}{Shora dolů vs. zdola nahoru}
        \begin{itemize}
            \visible<2->{\item Obě cesty řeší \textbf{tutéž množinu podproblémů} -- liší se jen \markdarkred{pořadím}, v~jakém je řeší.}
            \visible<3->{\item Při postupu zdola nahoru máme v~okamžiku, kdy nějaký podproblém řešíme, \textbf{už spočítané všechny menší}, na kterých závisí.}
            \visible<4->{\item Nepotřebujeme proto vůbec rekurzi -- stačí obyčejný cyklus a pole.}
        \end{itemize}
        \visible<5->{\begin{tnoticeframe}
            \begin{center}
                Výpočet zdola nahoru bývá i~o~něco rychlejší -- odpadá režie volání funkcí.
            \end{center}
        \end{tnoticeframe}}
    \end{frame}

    \subsection{Kešování}

    \begin{frame}{Kešování}
        \begin{itemize}
            \visible<2->{\item Mezi obyčejnou rekurzí a výpočtem zdola nahoru stojí \markdarkred{kešování} (angl. \emph{memoization}).}
            \visible<3->{\item Rekurzi ponecháme, jen si pořídíme tabulku $T$ již spočítaných hodnot -- \markdarkred{keš}.}
            \visible<4->{\item Na začátku každého volání se do keše podíváme:}
            \begin{itemize}
                \visible<5->{\item je-li hodnota vyplněná, \textbf{rovnou ji vrátíme},}
                \visible<6->{\item jinak ji poctivě spočítáme a \textbf{hned uložíme}.}
            \end{itemize}
        \end{itemize}
        \visible<7->{\obrazek[0.44\textheight]{images/kesovani.tex}}
    \end{frame}

    \begin{frame}{Kešování}
        \visible<2->{\begin{timportantframe}
            \begin{center}
                Kešování \textbf{prořezává strom rekurze}: opakovaná volání se z~něj stanou listy.
            \end{center}
        \end{timportantframe}}
        \begin{itemize}
            \visible<3->{\item Kešovaná rekurze je pořád \textbf{shora dolů}, ale už neplýtvá časem.}
            \visible<4->{\item Výpočet zdola nahoru dostaneme tak, že si uvědomíme, v~jakém pořadí se políčka tabulky vyplňují, a~toto pořadí zapíšeme rovnou cyklem.}
        \end{itemize}
    \end{frame}

    % =================================================================
    \section{Fibonacciho čísla}
    % =================================================================

    \titleframe{Fibonacciho čísla}

    \subsection{Rekurzivní výpočet}

    \begin{frame}{Fibonacciho čísla}
        \visible<2->{\begin{defblock}[Fibonacciho čísla]
            Posloupnost $(F_n)_{n=0}^{\infty}$ je dána předpisem
            \[F_0=0,\qquad F_1=1,\qquad F_n=F_{n-1}+F_{n-2}\ \text{ pro } n\geqslant 2.\]
        \end{defblock}}
        \visible<3->{\begin{center}
            $0,\ 1,\ 1,\ 2,\ 3,\ 5,\ 8,\ 13,\ 21,\ 34,\ 55,\ 89,\dots$
        \end{center}}
        \visible<4->{\begin{tinfoframe}
            \begin{center}
                Definice je rekurzivní $\implies$ nabízí se ji rovnou přepsat do algoritmu.
            \end{center}
        \end{tinfoframe}}
    \end{frame}

    \begin{frame}{Rekurzivní algoritmus}
        \begin{center}
            \begin{minipage}{0.85\textwidth}
                \begin{algorithm}[H]
                    \caption{\textsc{Fib}}
                    \KwIn{Číslo $n\in\N$}
                    \lIf{$n\leqslant 1$}{\Return{$n$}}
                    \Return{$\textsc{Fib}(n-1)+\textsc{Fib}(n-2)$}
                \end{algorithm}
            \end{minipage}
        \end{center}
        \visible<2->{\begin{tnoticeframe}
            \begin{center}
                Algoritmus je zjevně správný -- je to doslovný přepis definice. \textbf{Jak je ale rychlý?}
            \end{center}
        \end{tnoticeframe}}
    \end{frame}

    \begin{frame}{Strom rekurze pro $F_5$}
        \obrazek[0.85\textheight]{images/fib-strom.tex}
    \end{frame}

    \subsection{Složitost rekurzivního výpočtu}

    \begin{frame}{Složitost rekurzivního výpočtu}
        \begin{itemize}
            \visible<2->{\item V~listech stromu rekurze stojí samé jedničky a nuly.}
            \visible<3->{\item Hodnota každého vnitřního vrcholu je součtem hodnot listů pod ním $\implies$ abychom nasčítali $F_n$, musí být ve stromu \textbf{alespoň $F_n$ listů}.}
            \visible<4->{\item Přitom platí
            \[F_n\approx\frac{\varphi^{n}}{\sqrt{5}},\qquad \varphi=\frac{1+\sqrt{5}}{2}\approx 1{,}618,\]
            takže listů je \markdarkred{exponenciálně mnoho}.}
            \visible<5->{\item Časová složitost tedy činí $\bigO(\varphi^{n})$.}
        \end{itemize}
    \end{frame}

    \begin{frame}{Proč je to tak pomalé?}
        \obrazek[0.58\textheight]{images/fib-strom-opakovani.tex}
        \visible<2->{\begin{talertframe}
            \begin{center}
                Funkci $\textsc{Fib}$ voláme jen pro argumenty $0,1,\dots,n$, přesto ji voláme \textbf{exponenciálně mnohokrát}!
            \end{center}
        \end{talertframe}}
    \end{frame}

    \subsection{Kešovaný výpočet}

    \begin{frame}{Kešovaný výpočet}
        \begin{center}
            \begin{minipage}{0.92\textwidth}
                \begin{algorithm}[H]
                    \caption{\textsc{FibKes}}
                    \KwIn{Číslo $n\in\N$ a~tabulka $T$ (na počátku prázdná)}
                    \lIf{\textup{$T[n]$ je vyplněné}}{\Return{$T[n]$}}
                    \eIf{$n\leqslant 1$}{
                        $T[n]\gets n$\;
                    }{
                        $T[n]\gets\textsc{FibKes}(n-1)+\textsc{FibKes}(n-2)$\;
                    }
                    \Return{$T[n]$}
                \end{algorithm}
            \end{minipage}
        \end{center}
        \visible<2->{\begin{tnoticeframe}
            \begin{center}
                Rekurze zůstala, přibyl jen pohled do tabulky a~zápis do ní.
            \end{center}
        \end{tnoticeframe}}
    \end{frame}

    \begin{frame}{Prořezaný strom rekurze}
        \obrazek[0.85\textheight]{images/fib-strom-kes.tex}
    \end{frame}

    \begin{frame}{Složitost kešovaného výpočtu}
        \begin{itemize}
            \visible<2->{\item K~rekurzivnímu volání dojde jen tehdy, vyplňujeme-li políčko, v~němž dosud nic nebylo.}
            \visible<3->{\item Políček je $n+1$, takže vnitřních vrcholů stromu je nejvýše $n+1$.}
            \visible<4->{\item Pod každým z~nich leží nejvýše dva listy $\implies$ strom má nejvýše $3(n+1)$ vrcholů.}
            \visible<5->{\item V~každém vrcholu strávíme konstantní čas, celkem tedy $\bigO(n)$.}
        \end{itemize}
        \visible<6->{\begin{tnoticeframe}
            \begin{center}
                Z~$\bigO(\varphi^{n})$ jsme se dostali na $\bigO(n)$ -- a to jen díky tabulce.
            \end{center}
        \end{tnoticeframe}}
    \end{frame}

    \subsection{Výpočet zdola nahoru}

    \begin{frame}{Výpočet zdola nahoru}
        \begin{itemize}
            \visible<2->{\item K~výpočtu $T[k]$ potřebujeme jen $T[k-1]$ a~$T[k-2]$.}
            \visible<3->{\item Stačí tedy tabulku plnit v~pořadí $T[0],T[1],T[2],\dots$ a~vždy budeme mít vše, co potřebujeme.}
            \visible<4->{\item Rekurzi tím zcela odstraníme.}
        \end{itemize}
        \visible<5->{\obrazek[0.44\textheight]{images/fib-tabulka.tex}}
    \end{frame}

    \begin{frame}{Algoritmus zdola nahoru}
        \begin{center}
            \begin{minipage}{0.85\textwidth}
                \begin{algorithm}[H]
                    \caption{\textsc{FibTab}}
                    \KwIn{Číslo $n\in\N$}
                    $T[0]\gets 0,\ T[1]\gets 1$\;
                    \For{$k=2,3,\dots,n$}{
                        $T[k]\gets T[k-1]+T[k-2]$\;
                    }
                    \Return{$T[n]$}
                \end{algorithm}
            \end{minipage}
        \end{center}
        \begin{itemize}
            \visible<2->{\item Cyklus proběhne $(n-1)$-krát, každá iterace trvá konstantní čas $\implies$ $\bigO(n)$.}
            \visible<3->{\item Prostorová složitost je rovněž $\bigO(n)$; pamatujeme-li si jen poslední dvě hodnoty, klesne dokonce na $\bigO(1)$.}
        \end{itemize}
    \end{frame}

    \subsection{Porovnání výsledků}

    \begin{frame}{Porovnání výsledků}
        \visible<2->{\begin{center}
            \renewcommand{\arraystretch}{1.25}
            \begin{tabular}{|r|r|r|}
                \hline
                \rowcolor{decproblembg}
                $n$  & \textsc{Fib} (počet volání) & \textsc{FibTab} (počet sečtení) \\ \hline
                $10$ & $177$                       & $9$   \\ \hline
                $20$ & $21\,891$                   & $19$  \\ \hline
                $30$ & $2\,692\,537$               & $29$  \\ \hline
                $40$ & $331\,160\,281$             & $39$  \\ \hline
                $50$ & $40\,730\,022\,147$         & $49$  \\ \hline
            \end{tabular}
        \end{center}}
        \visible<3->{\begin{talertframe}
            \begin{center}
                Pro $n=50$ je rozdíl \textbf{devíti řádů} -- zlomek sekundy proti hodinám výpočtu.
            \end{center}
        \end{talertframe}}
    \end{frame}

    \begin{frame}{Porovnání výsledků}
        \obrazek[0.85\textheight]{images/slozitost-porovnani.tex}
    \end{frame}

    \subsection{Proč ten rozdíl vzniká?}

    \begin{frame}{Proč ten rozdíl vzniká?}
        \begin{itemize}
            \visible<2->{\item Různých podproblémů je jen $n+1$: hodnoty $F_0,F_1,\dots,F_n$.}
            \visible<3->{\item Strom rekurze je ale \textbf{exponenciálně velký}, protože tytéž podproblémy v~něm leží mnohokrát vedle sebe.}
            \visible<4->{\item Slepíme-li všechny vrcholy se stejným popiskem do jediného, dostaneme \markdarkred{graf podproblémů}.}
        \end{itemize}
        \visible<5->{\obrazek[0.44\textheight]{images/fib-graf-podproblemu.tex}}
    \end{frame}

    \begin{frame}{Proč ten rozdíl vzniká?}
        \visible<2->{\begin{timportantframe}
            \begin{center}
                Dynamické programování počítá s~\textbf{grafem podproblémů}, obyčejná rekurze se \textbf{stromem rekurze}.
            \end{center}
        \end{timportantframe}}
        \begin{itemize}
            \visible<3->{\item Graf podproblémů má $n+1$ vrcholů a~$\bigO(n)$ hran, kdežto strom rekurze $\bigO(\varphi^{n})$ vrcholů.}
            \visible<4->{\item Doba běhu dynamického programování je řádově \textbf{počet vrcholů plus počet hran} grafu podproblémů.}
            \visible<5->{\item Kešovaná rekurze je vlastně \textbf{prohledávání grafu podproblémů do hloubky}, výpočet zdola nahoru jeho průchod v~\textbf{obráceném topologickém pořadí}.}
        \end{itemize}
    \end{frame}

    % =================================================================
    \section{Dělení tyče}
    % =================================================================

    \titleframe{Dělení tyče}

    \subsection{Zadání úlohy}

    \begin{frame}{Dělení tyče}
        \begin{itemize}
            \visible<2->{\item Firma kupuje dlouhé ocelové tyče, řeže je na kratší kusy a~ty prodává.}
            \visible<3->{\item Kus délky $i$ prodá za cenu $p_i$; řezání je zdarma. Délky jsou vždy celá čísla.}
            \visible<4->{\item Ceny nemusí být úměrné délce -- proto se vyplatí přemýšlet, kde řezat.}
        \end{itemize}
        \visible<5->{\tyccenik}
        \visible<6->{\begin{tinfoframe}
            \begin{center}
                Anglicky se úloze říká \emph{rod cutting}.
            \end{center}
        \end{tinfoframe}}
    \end{frame}

    \begin{frame}{Formální zadání}
        \visible<2->{\decproblem{Dělení tyče}{
            Ceník $p_1,p_2,\dots,p_n$ a~délka tyče $n\in\N$.
        }{
            Jaký nejvyšší výnos $r_n$ lze získat rozřezáním tyče délky $n$ a~prodejem všech vzniklých kusů?
        }}
        \visible<3->{\begin{tnoticeframe}
            \begin{center}
                Tyč délky $n$ lze rozřezat $2^{n-1}$ způsoby -- na každém z~$n-1$ vnitřních míst se nezávisle rozhodujeme, zda řezat.
            \end{center}
        \end{tnoticeframe}}
    \end{frame}

    \begin{frame}{Všech osm způsobů pro $n=4$}
        \obrazek[0.62\textheight]{images/tyc-rezy.tex}
        \visible<2->{\begin{tnoticeframe}
            \begin{center}
                Nejlépe vyjde $4=2+2$ s~výnosem $r_4=10$, přestože celá tyč se prodá jen za $9$.
            \end{center}
        \end{tnoticeframe}}
    \end{frame}

    \subsection{Rekurzivní řešení}

    \begin{frame}{Optimální podstruktura}
        \begin{itemize}
            \visible<2->{\item Podívejme se na tyč jako na \textbf{první díl délky $i$} a~\textbf{zbytek délky $n-i$}.}
            \visible<3->{\item Zbytek je opět úloha o~dělení tyče, jen menší $\implies$ vyplatí se jej rozřezat \textbf{optimálně}.}
            \visible<4->{\item Kdyby totiž šel zbytek rozřezat lépe, zlepšili bychom i~řešení celé úlohy.}
        \end{itemize}
        \visible<5->{\obrazek[0.44\textheight]{images/tyc-rozklad.tex}}
    \end{frame}

    \begin{frame}{Rekurentní vztah}
        \visible<2->{\begin{propblock}
            Pro nejvyšší dosažitelný výnos $r_n$ platí
            \[r_0=0,\qquad r_n=\max_{1\leqslant i\leqslant n}\left(p_i+r_{n-i}\right)\ \text{ pro } n\geqslant 1.\]
        \end{propblock}}
        \begin{itemize}
            \visible<3->{\item Volba $i=n$ odpovídá tomu, že tyč vůbec neřežeme.}
            \visible<4->{\item Protože předem nevíme, které $i$ je nejlepší, \textbf{vyzkoušíme všechna}.}
        \end{itemize}
        \visible<5->{\begin{tinfoframe}
            \begin{center}
                Toto je přesně ta \textbf{optimální podstruktura}, o~které jsme mluvili.
            \end{center}
        \end{tinfoframe}}
    \end{frame}

    \begin{frame}{Rekurzivní algoritmus}
        \begin{center}
            \begin{minipage}{0.92\textwidth}
                \begin{algorithm}[H]
                    \caption{\textsc{DelTyc}}
                    \KwIn{Ceník $p[1\dots n]$ a~délka tyče $n$}
                    \lIf{$n=0$}{\Return{$0$}}
                    $q\gets-\infty$\;
                    \For{$i=1,2,\dots,n$}{
                        $q\gets\max\set{q,\ p[i]+\textsc{DelTyc}(p,n-i)}$\;
                    }
                    \Return{$q$}
                \end{algorithm}
            \end{minipage}
        \end{center}
        \visible<2->{\begin{talertframe}
            \begin{center}
                Algoritmus je správný, ale pro $n=40$ by počítal řádově hodiny.
            \end{center}
        \end{talertframe}}
    \end{frame}

    \begin{frame}{Strom rekurze pro $n=4$}
        \obrazek[0.85\textheight]{images/tyc-strom.tex}
    \end{frame}

    \begin{frame}{Složitost rekurzivního řešení}
        \begin{itemize}
            \visible<2->{\item Označme $T(n)$ počet volání funkce $\textsc{DelTyc}$ pro tyč délky $n$.}
            \visible<3->{\item Volání pro $n$ zavolá sebe sama pro všechny menší délky $0,1,\dots,n-1$, tedy
            \[T(0)=1,\qquad T(n)=1+\sum_{j=0}^{n-1}T(j).\]}
            \visible<4->{\item Odtud indukcí $T(n)=2^{n}$.}
            \visible<5->{\item Není divu: algoritmus opravdu probere \textbf{všech $2^{n-1}$ způsobů} rozřezání.}
        \end{itemize}
    \end{frame}

    \subsection{Dynamické programování}

    \begin{frame}{Graf podproblémů}
        \begin{itemize}
            \visible<2->{\item Různých podproblémů je přitom jen $n+1$: délky $0,1,\dots,n$.}
            \visible<3->{\item Slepením vrcholů stromu se stejným popiskem dostaneme opět \markdarkred{graf podproblémů}.}
        \end{itemize}
        \visible<4->{\obrazek[0.55\textheight]{images/tyc-graf-podproblemu.tex}}
    \end{frame}

    \begin{frame}{Rekurze s keší}
        \begin{center}
            \begin{minipage}{0.94\textwidth}
                \begin{algorithm}[H]
                    \caption{\textsc{DelTycKes}}
                    \KwIn{Ceník $p[1\dots n]$, délka $n$, prázdná tabulka $r$}
                    \lIf{\textup{$r[n]$ je vyplněné}}{\Return{$r[n]$}}
                    \eIf{$n=0$}{
                        $q\gets 0$\;
                    }{
                        $q\gets-\infty$\;
                        \For{$i=1,2,\dots,n$}{
                            $q\gets\max\set{q,\ p[i]+\textsc{DelTycKes}(p,n-i)}$\;
                        }
                    }
                    $r[n]\gets q$\;
                    \Return{$q$}
                \end{algorithm}
            \end{minipage}
        \end{center}
        \visible<2->{\begin{tnoticeframe}
            \begin{center}
                Stejný trik jako u~Fibonacciho čísel -- rekurze zůstává, jen si výsledky pamatujeme.
            \end{center}
        \end{tnoticeframe}}
    \end{frame}

    \begin{frame}{Výpočet zdola nahoru}
        \begin{itemize}
            \visible<2->{\item Podproblém délky $j$ závisí jen na kratších tyčích $0,1,\dots,j-1$.}
            \visible<3->{\item Stačí tedy políčka $r[0],r[1],\dots,r[n]$ vyplňovat \textbf{v pořadí rostoucí délky}.}
        \end{itemize}
        \visible<4->{\obrazek[0.50\textheight]{images/tyc-tabulka.tex}}
    \end{frame}

    \begin{frame}{Algoritmus zdola nahoru}
        \begin{center}
            \begin{minipage}{0.92\textwidth}
                \begin{algorithm}[H]
                    \caption{\textsc{DelTycTab}}
                    \KwIn{Ceník $p[1\dots n]$ a~délka tyče $n$}
                    $r[0]\gets 0$\;
                    \For{$j=1,2,\dots,n$}{
                        $q\gets-\infty$\;
                        \For{$i=1,2,\dots,j$}{
                            $q\gets\max\set{q,\ p[i]+r[j-i]}$\;
                        }
                        $r[j]\gets q$\;
                    }
                    \Return{$r[n]$}
                \end{algorithm}
            \end{minipage}
        \end{center}
        \visible<2->{\begin{tnoticeframe}
            \begin{center}
                Rekurze je pryč -- zbyly dva vnořené cykly a~jedno pole.
            \end{center}
        \end{tnoticeframe}}
    \end{frame}

    \begin{frame}{Výsledná tabulka pro náš ceník}
        \visible<2->{\tyccenik}
        \visible<3->{\begin{center}
            \renewcommand{\arraystretch}{1.25}
            \begin{tabular}{|r|c|c|c|c|c|c|c|c|c|c|c|}
                \hline
                \rowcolor{decproblembg}
                $j$    & $0$ & $1$ & $2$ & $3$ & $4$ & $5$ & $6$ & $7$ & $8$ & $9$ & $10$ \\ \hline
                $r[j]$ & $0$ & $1$ & $5$ & $8$ & $10$ & $13$ & $17$ & $18$ & $22$ & $25$ & $30$ \\ \hline
            \end{tabular}
        \end{center}}
        \visible<4->{\begin{tinfoframe}
            \begin{center}
                Např. $r_7=18$ odpovídá rozřezání $7=1+6$, kdežto $r_{10}=30$ znamená tyč vůbec neřezat.
            \end{center}
        \end{tinfoframe}}
    \end{frame}

    \subsection{Porovnání složitostí}

    \begin{frame}{Porovnání složitostí}
        \begin{itemize}
            \visible<2->{\item Vnější cyklus proběhne $n$-krát, vnitřní pro dané $j$ celkem $j$-krát.}
            \visible<3->{\item Dohromady tedy
            \[\sum_{j=1}^{n}j=\frac{n(n+1)}{2}=\bigO(n^{2})\]
            iterací, každá v~konstantním čase.}
            \visible<4->{\item Prostorová složitost činí $\bigO(n)$ -- jedno pole délky $n+1$.}
        \end{itemize}
        \visible<5->{\begin{tinfoframe}
            \begin{center}
                Kešovaná rekurze má tutéž časovou složitost $\bigO(n^{2})$ -- řeší přesně tytéž podproblémy, jen v~jiném pořadí.
            \end{center}
        \end{tinfoframe}}
    \end{frame}

    \begin{frame}{Porovnání složitostí}
        \visible<2->{\begin{center}
            \renewcommand{\arraystretch}{1.35}
            \begin{tabular}{|l|c|c|}
                \hline
                \rowcolor{decproblembg}
                                  & čas         & paměť     \\ \hline
                $\textsc{DelTyc}$ (rekurze) & $\bigO(2^{n})$ & $\bigO(n)$ \\ \hline
                $\textsc{DelTyc}$ s~keší    & $\bigO(n^{2})$ & $\bigO(n)$ \\ \hline
                $\textsc{DelTycTab}$ (zdola nahoru) & $\bigO(n^{2})$ & $\bigO(n)$ \\ \hline
            \end{tabular}
        \end{center}}
        \visible<3->{\begin{timportantframe}
            \begin{center}
                Z~exponenciálního algoritmu se stal \textbf{kvadratický}, aniž bychom potřebovali víc paměti.
            \end{center}
        \end{timportantframe}}
    \end{frame}

    \subsection{Rekonstrukce řešení}

    \begin{frame}{Jak zjistit, kde řezat?}
        \begin{itemize}
            \visible<2->{\item Zatím umíme spočítat jen \textbf{výnos}, nikoli samotné řezy.}
            \visible<3->{\item Stačí si při vyplňování tabulky navíc pamatovat pole $s$, kde $s[j]$ je délka \textbf{prvního dílu} optimálního rozřezání tyče délky $j$.}
            \visible<4->{\item Řešení pak vypíšeme postupně: vytiskneme $s[n]$, zkrátíme tyč na $n-s[n]$ a~opakujeme, dokud něco zbývá.}
        \end{itemize}
        \visible<5->{\begin{tinfoframe}
            \begin{center}
                Je to táž myšlenka jako uchovávání předchůdců při hledání nejkratší cesty v~grafu.
            \end{center}
        \end{tinfoframe}}
    \end{frame}

    % =================================================================
    \section{Shrnutí}
    % =================================================================

    \begin{frame}{Princip dynamického programování}
        \begin{enumerate}
            \visible<2->{\item Začneme \textbf{rekurzivním algoritmem}, který je exponenciálně pomalý.}
            \visible<3->{\item Odhalíme \textbf{opakované výpočty} týchž podproblémů.}
            \visible<4->{\item Pořídíme si \textbf{tabulku} (keš) a~pamatujeme si v~ní, co už jsme spočítali. Tím prořežeme strom rekurze.}
            \visible<5->{\item Uvědomíme si, v~jakém \textbf{pořadí} lze tabulku vyplňovat, a~rekurzi nahradíme cyklem \textbf{zdola nahoru}.}
        \end{enumerate}
        \visible<6->{\begin{timportantframe}
            \begin{center}
                Klíčem je vždy \textbf{graf podproblémů}: kolik má vrcholů a~hran, tak dlouho výpočet trvá.
            \end{center}
        \end{timportantframe}}
    \end{frame}

    \zdrojeframe{mares-valla-labyrint, cormen-algorithms}

\end{document}
